\scanpage{5}
degrés positifs. On voit que si $i$ est effectif, il en est de \unsure{$\overline{\mathrm m}$|même} de ses conjugués par éléments de $G(C)$, et de ses conjugués par autom. de $C$, en particulier de $\bar i$. \deleted{\unsure{De}}D'autre part, pour $i$ donné, la sous-catégorie \inserted{pleine} $P(i)$ de $M$ \unsure{formé|formée} des $E$ tels que $E_C$ soit à $i$-graduation positive est \inserted{str. pleine,} stable par $\deleted{\otimes}$ sommes, par \deleted{\unsure{les}} sous-objet, par $\otimes$, contient \inserted{$0$, et} $\Lambda(-1)$ si on suppose que $\deleted{\varepsilon i(\lambda)=\lambda^\nu\ \text{avec }\nu}\deleted{\langle i,\varepsilon\rangle}\inserted{\varepsilon(i)}\geq 0$, et satisfait la condition que $\forall E\in\operatorname{Ob}M$, $\exists n$ tel que $E(\unsure{-n|n})\in\operatorname{Ob}\deleted{\unsure{M}}P(i)$, si on suppose $\boxed{\deleted{\langle i,\varepsilon\rangle}\inserted{\varepsilon(i)}>0}$ [N.B. on désigne par $\langle i,\varepsilon\rangle$ l'entier $\nu$ tel que $\varepsilon(i(\lambda))=\lambda^\nu$]. \deleted{Si l'on pose \unsure{égalem.t|engendré} pour les conjugués $\operatorname{int}(g)\circ i$ et les conjugués $\,{}^{\sigma}i$ $(g\in G(C),\ i\in$} \inserted{Prenant des intersections} des sous-catégories $P(i)$ on trouve des $M^+$ satisfaisant aux axiomes considérés, sauf peut être au dernier sur la graduation. \inserted{Par intersection de tels \unsure{jouets|fonct's|foncteurs}, itou \ldots}

Les $i$ pour lesquels, $M^+$ étant donné,
% archive sheet number: 4
